\documentclass[11pt]{beamer}

\usepackage[utf8]{inputenc}   
\usepackage[T1]{fontenc}     
\usepackage[polish]{babel}   

\graphicspath{{Images/}{./}} 
\usepackage{booktabs} 
\usetheme{Warsaw}
\useinnertheme{circles}
\usepackage{palatino} 
\usepackage[default]{opensans} 




\title[Automatyzacja procesów biznesowych z IDP]{Automatyzacja procesów biznesowych z wykorzystaniem inteligentnego przetwarzania formularzy IDP}

\author[Demski, Dziubich, Formela]{
    Paweł Demski \and 
    Paweł Dziubich \and 
    Mateusz Formela
}

\institute[PG]{
    Promotor: dr inż. Krzysztof Bikonis\\
    \smallskip 
    \textit{Ktedra Systemów Geoinformatycznych}
}

\date{\today}

%----------------------------------------------------------------------------------------

\begin{document}

\begin{frame}
    \titlepage 
\end{frame}

% \begin{frame}
%     \frametitle{Plan prezentacji}
%     \tableofcontents
% \end{frame}

%----------------------------------------------------------------------------------------
\section{Cel i Założenia}
%----------------------------------------------------------------------------------------


\begin{frame}
    \frametitle{Cel pracy – Kontekst: Technologia IDP}
    \textbf{Inteligentne Przetwarzanie Dokumentów (IDP)} to technologia AI automatyzująca analizę i ekstrakcję danych z dokumentów (faktury, umowy, skany).
    
    \bigskip
    \begin{itemize}
        \item \textbf{Fundamenty:} Wykorzystanie OCR (Optical Character Recognition), NLP (Natural Language Processing ) oraz ML (Machine learning).
        \item \textbf{Działanie:} Przekształcanie nieustrukturyzowanego chaosu w uporządkowane dane gotowe dla systemów ERP/CRM.
        \item \textbf{Korzyści:} Zrozumienie kontekstu (nawet notatek odręcznych), usprawnienie procesów biznesowych i redukcja pracy ręcznej.
    \end{itemize}
\end{frame}


\begin{frame}
    \frametitle{Cel pracy – Zadania projektowe}
    Głównym celem pracy jest \textbf{wdrożenie chmurowe oprogramowania BPM}(Business Process Management) oraz jego adaptacja do wybranej dziedziny (np. administracji).
    
    \bigskip
    \begin{block}{Zakres realizacji}
        \begin{itemize}
            \item \textbf{Modelowanie:} Zaprojektowanie procesu biznesowego opartego na IDP.
            \item \textbf{Integracja (V)LLM:} Wybór i wdrożenie modelu klasy Vision Large Language Model do przetwarzania danych niestrukturyzowanych.
            \item \textbf{Ewaluacja:} Ocena jakościowa i wydajnościowa stworzonego rozwiązania.
        \end{enumerate}
    \end{block}
\end{frame}

%----------------------------------------------------------------------------------------
\section{Koncepcja Rozwiązania}
%----------------------------------------------------------------------------------------

\begin{frame}
    \frametitle{Koncepcja CloudIDP Flow}
    Projekt to chmurowa platforma automatyzacji procesów biznesowych (BPM) zintegrowana z modułem Inteligentnego Przetwarzania Dokumentów (IDP).

    \bigskip
    \begin{itemize}
        \item \textbf{Adresowany problem:} Ręczne, czasochłonne i podatne na błędy wprowadzanie danych z nieustrukturyzowanych dokumentów.
        \item \textbf{Podejście Low-code/No-code:} Intuicyjne środowisko graficzne pozwalające analitykom biznesowym na samodzielne modelowanie procesów bez umiejętności programistycznych.
        \item \textbf{Automatyzacja decyzji:} Wykorzystanie dużych modeli językowych do podejmowania prostych decyzji na podstawie wyekstrahowanych danych.
    \end{itemize}
\end{frame}

\begin{frame}
    \frametitle{Użytkownicy i ich role}
    Platforma definiuje dwa główne typy interakcji:
    
    \begin{columns}[t]
        \begin{column}{0.45\textwidth}
            \textbf{Analitycy Biznesowi}
            \begin{itemize}
                \item Projektowanie przepływów pracy (Citizen Developers).
                \item Konfiguracja modułu IDP i mapowanie danych do systemów ERP/CRM.
            \end{itemize}
        \end{column}
        \begin{column}{0.45\textwidth}
            \textbf{Pracownicy Administracji}
            \begin{itemize}
                \item Nadzór nad procesem i zatwierdzanie danych w przypadku niskiej pewności modelu.
                \item Korzystanie z ustrukturyzowanych informacji w systemach docelowych.
            \end{itemize}
        \end{column}
    \end{columns}
\end{frame}

%----------------------------------------------------------------------------------------
\section{Architektura Systemu}
%----------------------------------------------------------------------------------------

\begin{frame}
    \frametitle{Architektura Systemu}
    \begin{itemize}
        \item \textbf{Środowisko chmurowe:} Skalowalne rozwiązanie BPM wdrożone w infrastrukturze chmurowej (np. AWS, Azure).
        \item \textbf{Integracja (V)LLM:} Wykorzystanie modeli Vision Large Language Model do wielwymiarowej analizy i ekstrakcji danych.
        \item \textbf{Ekosystem integracyjny:} 
        \begin{itemize}
            \item Współpraca z zewnętrznymi systemami ERP/CRM klienta.
            \item Wykorzystanie usług katalogowych do uwierzytelniania.
            \item Bezpieczeństwo danych zgodnie z wymogami RODO.
        \end{itemize}
    \end{itemize}
\end{frame}

%----------------------------------------------------------------------------------------
\section{Technologie}
%----------------------------------------------------------------------------------------

\begin{frame}
    \frametitle{Wykorzystywane Technologie}
    \begin{table}
        \begin{tabular}{l l}
            \toprule
            \textbf{Warstwa} & \textbf{Technologia / Narzędzie} \\
            \midrule
            Middleware / BPM & n8n / Camunda \\
            AI / IDP & Modele (V)LLM (OCR, NLP, ML)  \\
            Infrastruktura & AWS / Azure / TASKCloud  \\
            Integracje & Python / Node.js \\
            Współdzielenie kodu & Git / GitHub  \\
            CI/CD i Testy & GitHub Actions / Postman  \\
            \bottomrule
        \end{tabular}
    \end{table}
\end{frame}
%----------------------------------------------------------------------------------------
\section{Status i Harmonogram}
%----------------------------------------------------------------------------------------

\begin{frame}
    \frametitle{Aktualny stan prac}
    \begin{itemize}
        \item \textbf{Research:} 
        \item \textbf{Design:} 
        \item \textbf{Prototypowanie:} 
    \end{itemize}
\end{frame}

\begin{frame}
    \frametitle{Harmonogram prac}
    \small 
    \begin{itemize}
        \setlength\itemsep{0.8em} 
        \item \textbf{Marzec:} Zapoznanie się z dziedziną pracy
        \item \textbf{Kwiecień -- Maj:} Przegląd literatury przedmiotu i selekcja kluczowego procesu
        \item \textbf{Czerwiec -- Lipiec:} Przygotowanie środowiska do przeprowadzenia eksperymentów
        \item \textbf{Sierpień -- Wrzesień:} Wdrożenie pilotażowe zautomatyzowanego procesu opartego na IDP
        \item \textbf{Październik -- Grudzień:} Badanie oraz ocena wdrożonych usprawnień
        \item \textbf{Grudzień:} Redakcja pracy
    \end{itemize}
\end{frame}

%----------------------------------------------------------------------------------------

% \begin{frame}[plain]
%     \begin{center}
%         {\Huge Koniec}
%         \bigskip\bigskip
%         {\LARGE Pytania i komentarze?}
%     \end{center}
% \end{frame}

\end{document}